\section{積立金の積立て及び運用}

\subsection{積立金の積立てと財政検証}


\begin{InyouJobunbox}
  \begin{itemize}
    \item 法第59条〜第61条
  \end{itemize}
\end{InyouJobunbox}

\begin{fillblank}[自作 2026/11/9]
  \noindent \textbf{◯DB法}
  \begin{lawstructure}
    \begin{lawarticle}{第五十九条}{積立金の積立て}
      事業主等は、（　　A　　）において、給付に充てるべき積立金（以下「積立金」という。）を積み立てなければならない。
    \end{lawarticle}

    \begin{lawarticle}{第六十条}{積立金の額}
      積立金の額は、加入者及び加入者であった者（以下「加入者等」という。）に係る次項に規定する責任準備金の額及び第三項に規定する最低積立基準額を下回らない額でなければならない。

      \lawparagraph{2}
      責任準備金の額は、当該事業年度の末日における（　　B　　）から（　　C　　）を控除した額を基準として、厚生労働省令で定めるところにより算定した額とする。

      \lawparagraph{3}
      最低積立基準額は、加入者等の（　　D　　）として政令で定める基準に従い規約で定めるものに要する費用の額の予想額を計算し、これらの予想額の合計額の現価として厚生労働省令で定めるところにより算定した額とする。
    \end{lawarticle}

    \begin{lawarticle}{第六十一条}{決算における責任準備金の額等の計算}
      事業主等は、（　　E　　）において、積立金の額が前条第二項に規定する責任準備金の額（以下「責任準備金の額」という。）及び同条第三項に規定する最低積立基準額（以下「最低積立基準額」という。）を上回っているかどうかを計算しなければならない。
    \end{lawarticle}
  \end{lawstructure}
\end{fillblank}

\begin{answerbox}
  \begin{enumerate}[A: ]
    \item 毎事業年度の末日
    \item 給付に要する費用の額の予測額の現価
    \item 掛金収入の額の予想額の現価
    \item 当該事業年度の末日までの加入者期間に係る給付
    \item 毎事業年度の決算
  \end{enumerate}
\end{answerbox}

\newpage

\subsection{責任準備金}

\begin{InyouJobunbox}
  \begin{itemize}
    \item 法第60条
    \item 則第53条、第63条
  \end{itemize}
\end{InyouJobunbox}

\begin{fillblank}[数理人試験 H29 問題1 設問1]
  \noindent \textbf{◯DB則}
  \begin{lawstructure}
    \begin{lawarticle}{第五十三条}{責任準備金の額}
      責任準備金の額は、当該事業年度の末日における（　　A　　）の現価と
      （　　B　　）を合算した額から、掛金の額（標準掛金額及び補足掛金額を
      合算した額又は（　　C　　）をいう。第三項において同じ。）
      の現価に相当する額と（　　B　　）に対応するために追加的に拠出される
      こととなる掛金の額の予想額（同項において「（　　D　　）」という。）の現価に
      相当する額を合算した額を控除した額とする。
    
      \lawparagraph{2}
      前項の予想額の現価の計算は、（　　E　　）の基礎率を用いて行うものとする。

      \lawparagraph{3}
      （　　D　　）の現価に相当する額は、（　　B　　）から（　　F　　）
      （積立金の額と（　　G　　）の現価を合算した額から（　　A　　）
      の現価に相当する額を控除した額（当該額が零未満となる場合にあっては零とする。）
      をいう。）を控除した額（当該額が零未満となる場合にあっては零とする。）とする。
    \end{lawarticle}
  \end{lawstructure}
\end{fillblank}

\begin{answerbox}
  \begin{enumerate}[A: ]
    \item 通常予測給付額
    \item 財政悪化リスク相当額
    \item リスク分担型企業年金掛金額
    \item 追加拠出可能額
    \item 前回の財政計算
    \item リスク充足額
    \item 掛金の額の予想額
  \end{enumerate}
\end{answerbox}

\begin{truefalse}[年金2 H29 問題1(1)\textcircled{1}]
  下線部の正誤を判定し、誤っている場合は正しい内容に改めなさい。
\\

  確定給付企業年金法施行規則第 53 条において、責任準備金の額は、「通常予測給付額の現価
と\underline{リスク充足額}を合算した額から、掛金の額の現価に相当する額と追加拠出可能額の現価に
相当する額を合算した額を控除した額」とされている。
\end{truefalse}

\begin{answerbox}
  \textbf{誤}

  財政悪化リスク相当額
\end{answerbox}

\newpage

\subsection{継続基準の財政検証}

\begin{InyouJobunbox}
  \begin{itemize}
    \item 法第60条第1項、第62条
    \item 則第56条〜第57条、第63条
  \end{itemize}
\end{InyouJobunbox}

\begin{truefalse}[年金2 2025 問題1(1)ア\ajRoman{1}]
  以下の文章の正誤を判定し、誤っている場合は正しい内容に改めなさい。
\\

  事業主等は、積立金の額が責任準備金の額に照らし厚生労働省令で定めるところ
  により算定した額を下回っている場合には、掛金の額を再計算しなければならな
  い。この積立金の額は、掛金の額の計算に用いる積立金の額の評価方法を用いて計算する。
\end{truefalse}

\begin{answerbox}
  \textbf{正}
\end{answerbox}

\begin{shortanswer}[年金2 2024 問題2(1)]
  確定給付企業年金法施行規則第56条に定める許容繰越不足金の算定方法を3つ簡記しなさい。 
  なお、当該規則に記載の「規約で定める率」の範囲についても触れた上で解答を入力すること。
   (400 字以内)
\end{shortanswer}

\begin{answerbox}
  \begin{enumerate}[(1)]
    \item 当該事業年度以後２０年間における標準掛金額の予想額の現価×規約で定める率（規約で定める率は１００分の１５を超えないこと、また、標準掛金額の予想額の現価は掛金算定に用いた予定利率を用いて算定）
    \item 当該事業年度の末日における責任準備金の額×規約で定める率（規約で定める率は１００分の１５（数理的評価により積立金の額を評価する場合は、１００分の１０）を超えてはならない）
    \item (1)又は(2)のうちいずれか小さい額
  \end{enumerate}
\end{answerbox}

\newpage

\subsection{最低積立基準額}

\begin{InyouJobunbox}
  \begin{itemize}
    \item 法第60条〜第61条
    \item 令第37条
    \item 則第54条〜第55条、第63条
    \item 解釈第5-1
    \item 告示（則第55条第1項第1号に規定する予定利率）
  \end{itemize}
\end{InyouJobunbox}

\subsubsection{最低保全給付の計算方法}


\begin{fillblank}[年金2 2024 問題1(2)C〜E　\&　数理人試験 2021 問題1 設問2 A〜C]
  \noindent \textbf{◯DB則}
  \begin{lawstructure}
    \begin{lawarticle}{第五十四条}{（　　A　　）}
      令第三十七条第五号及び第六号に定める加入者が老齢給付金又は脱退一時金
      （法第四十一条第二項第一号に係るものに限る。以下この条において同じ。）を
      受けるための要件を満たした場合に支給されることとなる当該老齢給付金及び当該脱退一時金
      のうち当該加入者の当該事業年度の末日までの加入者期間に係る分の額は、
      次に掲げる方法又はこれらに準ずる方法により計算するものとする。

      \lawitem{一}
      当該加入者が加入者の資格を喪失する（　　B　　）に達した日において加入者の資格
      を喪失する場合に支給されることとなる老齢給付金の額又は脱退一時金の額に、
      加入者が加入者の資格を取得した日から当該（　　B　　）に達するまでの加入者期間
      のうち当該事業年度の末日までの加入者期間に係る分として定めた率を乗ずる方法

      \lawitem{二}
      当該事業年度の末日において当該加入者が加入者の資格を喪失した場合に支給されること
      となる老齢給付金の額（第二十七条第一号の（　　C　　）を行うこととなっている場合にあっては、
      当該（　　C　　）を行わないものとして計算した額）又は脱退一時金の額に（　　D　　）に応じて
      定めた率を乗ずる方法

      \lawparagraph{2}
      法第二十八条第三項の規定に基づく加入者となる前の期間の加入者期間への算入又は
      給付の額の増額（以下この項において「給付改善等」という。）を行う場合にあっては、
      令第三十七条各号に定める加入者等の当該事業年度の末日までの加入者期間に係る給付
      として規約で定めるもの（以下「最低保全給付」という。）の額は、当該給付改善等に
      より増加する給付の額に、当該給付改善等に係る規約が効力を有することとなる日から
      当該事業年度の末日までの年数（その期間に一年に満たない端数がある場合にあっては、
      これを切り捨てるものとする。）を（　　E　　）から減じた数（当該数が零未満となる場合にあっては、
      零とする。）を（　　E　　）で除して得た数を乗じて得た額を、前項の規定に基づき計算した額から控除
      した額とすることができる。
    \end{lawarticle}
  \end{lawstructure}
\end{fillblank}

\begin{answerbox}
  \begin{enumerate}[A: ]
    \item 最低保全給付の計算方法
    \item 標準的な年齢
    \item 加算
    \item 当該加入者の年齢
    \item 五
  \end{enumerate}
\end{answerbox}

\begin{fillblank}[数理人試験 H27 問題1 設問3]
  \noindent \textbf{◯DB令}
  \begin{lawstructure}
    \begin{lawarticle}{第三十七条}{過去の加入者期間に係る給付の基準}
      法第六十条第三項の政令で定める基準は、加入者等の当該事業年度の末日までの加入者期間に係る給付として（　　A　　）で定めるものが、次の各号に掲げる者の区分に応じ、当該各号に定めるものであることとする。

      \lawitem{一}
      当該事業年度の末日において、年金給付の支給を受けている者　当該年金給付

      \lawitem{二}
      当該事業年度の末日において、老齢給付金の受給権者であって法第三十七条第一項の規定に基づき
      その老齢給付金の支給の（　　B　　）の申出をしている者　その者が当該事業年度の末日において当該支給の（　　B　　）の申出をした老齢給付金の支給を請求するとした場合における年金として支給される老齢給付金

      \lawitem{三}
      当該事業年度の末日において、老齢給付金を受けるための要件のうち（　　C　　）以外の要件を満たす者（
      加入者及び老齢給付金の（　　D　　）に代えて脱退一時金の支給を受けた者を除く。）　その者が（　　C　　）を満たしたときに年金として支給される老齢給付金

      \lawitem{四}
      当該事業年度の末日において、法第四十一条第二項第一号に係る脱退一時金の受給権者であって、
      同条第四項の規定に基づきその脱退一時金の（　　E　　）の支給の（　　B　　）の申出をしている者　その者が当該事業年度の末日において、脱退一時金の支給を請求するとした場合に支給される脱退一時金

      \lawitem{五}
      当該事業年度の末日において、加入者であって、老齢給付金を受けるための要件のうち（　　C　　）以外の要件
      を満たす者　その者が老齢給付金を受けるための要件を満たしたときに支給される当該老齢給付金のうち、
      その者の当該事業年度の末日までの（　　F　　）として、厚生労働省令で定めるところにより計算した額


      \lawitem{六}
      当該事業年度の末日における加入者（前号に規定する者を除く。）　その者が脱退一時金を
      受けるための要件を満たしたときに支給される当該脱退一時金のうち、その者の当該事業年度の末日までの
      （　　F　　）として、厚生労働省令で定めるところにより計算した額

    \end{lawarticle}
  \end{lawstructure}
\end{fillblank}

\begin{answerbox}
  \begin{enumerate}[A: ]
    \item 規約
    \item 繰下げ
    \item 老齢給付金支給開始要件
    \item 全部
    \item 全部または一部
    \item 加入者期間に係る分
  \end{enumerate}
\end{answerbox}

\begin{fillblank}[数理人試験 2025 問題1 設問1 A〜D]
  \noindent \textbf{◯法令解釈}
  \begin{lawstructure}
    \begin{lawarticle}{}{}
      \noindent 第５ 積立金の積立てに関する事項

      １ 最低積立基準額の算定においては、当該額が確定給付企業年金が終了した場合における残余財産の分
        配額の算定基礎等となることを踏まえて、次の点に留意すること。
      
      \lawitem{(1)}
      令第３７条第１号の「当該年金給付」の取扱いは、次のとおりとすること。

      \hangedlawitem{\textcircled{1}}{3}
      あらかじめ給付の額が（　　A　　）することとなっている場合にあっては、当該（　　A　　）
を織り込むこと。
      \hangedlawitem{\textcircled{2}}{3}
    支給停止中の給付については、停止終了後に支給されることが確実な場合にあっては、当該給付
を計算に織り込むこと（例えば、障害給付金が支給されているために老齢給付金が支給停止と
なっている場合であって、当該障害給付金の支給期間が終了した場合に老齢給付金が支給され
ることとなっているときは、当該老齢給付金を織り込むこと。）。

      \lawitem{(2)}
規則第５４条第１項の「これらに準ずる方法」とは、同項第１号及び第２号の方法を組み合わせた
方法をいうこと。

      \lawitem{(3)}
規則第５４条第１項第２号の「加入者の（　　B　　）に応じて定めた率」は、 基本的に、
（　　C　　）又は（　　D　　）までの期間に応じて定める率であるが、この場合において、
（略）。

    \end{lawarticle}
  \end{lawstructure}
\end{fillblank}

\begin{answerbox}
  \begin{enumerate}[A: ]
    \item 増加又は減少
    \item 年齢
    \item 一定率
    \item 支給開始時
  \end{enumerate}
\end{answerbox}


\begin{shortanswer}[年金2 2019 問題2(1)\textcircled{1}]
  最低積立基準額の算定に関して、確定給付企業年金法施行規則第 54 条第 1 項第 2 号に規定する「加入者の年齢に応じて定め
た率」を設定する場合における留意点について簡記しなさい。
\end{shortanswer}

\begin{answerbox}
  基本的に、一定率または支給開始時までの期間に応じて定めた率であるが、
この場合において、加入者の年齢が増すごとに減少するものではないこと
\end{answerbox}

\newpage

\subsubsection{最低積立基準額（リスク分担型以外の場合）}



\begin{fillblank}[数理人試験 2025 問題1 設問1 E〜G]
  \noindent \textbf{◯法令解釈}
  \begin{lawstructure}
    \begin{lawarticle}{}{}
      \noindent 第５ 積立金の積立てに関する事項

      １ 最低積立基準額の算定においては、当該額が確定給付企業年金が終了した場合における残余財産の分
        配額の算定基礎等となることを踏まえて、次の点に留意すること。
      
      \lawitem{(4)}
      確定給付企業年金法施行規則第５５条第１項第１号に規定する予定利率（平成１４年厚生労働省告
      示第５９号）に規定する「（　　A　　）以内の率」を設定するときは、その設定の根拠及び最低
      積立基準額に及ぼす影響について、（　　B　　）において十分な検討を行っている必要があり、
      （　　C　　）に対して十分な情報提供を行うこと。 なお、「（　　A　　）以内の率」を設定し
      ている場合に確定給付企業年金の終了、確定拠出年金（企業型）への移行又は中小企業退職金共済
      への移換をするときは、最低積立基準額の算定に用いる予定利率について改めて（　　B　　）に
      おいて十分な検討を行い、検討の結果として「（　　A　　）以内の率」を設定するときは、
      （　　C　　） に対して十分な情報提供を行うこと。

    \end{lawarticle}
  \end{lawstructure}
\end{fillblank}

\begin{answerbox}
  \begin{enumerate}[A: ]
    \item 0.5パーセント
    \item 労使間や代議員会
    \item 加入者及び受給権者等
  \end{enumerate}
\end{answerbox}

\begin{truefalse}[年金2 2023 問題1(1)イ]
  下線部の正誤を判定し、誤っている場合は正しい内容に改めなさい。
\\

  確定給付企業年金において、最低積立基準額の算定に用いる予定死亡率は、基準死亡率
に、\underline{加入者等が男子である}\\
\underline{場合にあっては 0.72 を、加入者等が女子である場合にあっては
0.72}を、それぞれ乗じて得た率と定められている。
\end{truefalse}

\begin{answerbox}
  \textbf{誤}

  加入者等が男子である場合にあっては、0.84を、
  加入者等が女子である場合にあっては、0.825
\end{answerbox}

\begin{explanationbox}
  ＜類題＞

  数理人試験 2021 問題1 設問2 D,E

  年金２ 2025 問題1(1)ウ\ajRoman{2}
\end{explanationbox}

\begin{fillblank}[自作 2026/11/10]
  最低積立基準額の計算の基礎となる予定利率は、令和 8 年度は年率 （　　A　　）\%(当該年率に 0.5\%以内の率を加減して得た率を予定利率とすることについて、労働組合等の同意を得た場合にあっては当該加減して得た率)である。
\end{fillblank}

\begin{answerbox}
  1.64
\end{answerbox}

\begin{shortanswer}[年金2 2019 問題2(1)\textcircled{2}]
  最低積立基準額の算定に関して、
  確定給付企業年金法施行規則第 55 条第1 項第1 号に規定する予定利率の「0.5 パーセント以
内の率」を設定する場合における留意点について簡記しなさい。
\end{shortanswer}

\begin{answerbox}
  その設定の根拠及び最低積立基準額に及ぼす影響について、労使間や代議員会において十分な検討を行っていること必要があり、加入者等および受給権者等に対して十分な情報提供を行うこと。なお、「0.5パーセント以内の率」を設定している場合に確定給付企業年金の終了、確定拠出年金（企業型）への移行または中小企業退職金共済への移換をする時は、最低積立基準額の算定に用いる予定利率について改めて労使間や代議員会において十分な検討を行い、検討の結果として「0.5パーセント以内の率」を設定するときは、加入者及び受給権者等に対して十分な情報提供を行うこと。
\end{answerbox}

\newpage

\subsubsection{最低積立基準額（リスク分担型の場合）}


\begin{fillblank}[数理人試験 H29 問題1 設問1 h　\&　数理人試験 2021 問題1 設問2 F,G]
  \noindent \textbf{◯DB則}
  \begin{lawstructure}
    \begin{lawarticle}{第五十五条}{最低積立基準額}
      法第六十条第三項の厚生労働省令で定めるところにより算定した額の計算の基礎となる予定利率及び予定死亡率は、次のとおりとする。

      （略）

      \lawparagraph{3}
      （　　A　　）を実施している場合にあっては、法第六十条第三項の現価の算定において、（　　B　　）を第一項に規定する予定利率及び予定死亡率並びに前項に規定する指標の予測を算定の基礎とするならば算定されることとなる法第六十条第三項の現価（　　C　　）を計算の基礎とするものとする。
    \end{lawarticle}
  \end{lawstructure}
\end{fillblank}

\begin{answerbox}
  \begin{enumerate}[A: ]
    \item リスク分担型企業年金
    \item 積立金の額
    \item で除して得た率
  \end{enumerate}
\end{answerbox}

\newpage

\subsection{非継続基準の財政検証}


\begin{InyouJobunbox}
  \begin{itemize}
    \item 法第60条第1項、第63条
    \item 則第58条〜第59条、第63条、則附則(H24年第13号)第4条
    \item 解釈第5-2
  \end{itemize}
\end{InyouJobunbox}

\subsubsection{積立比率方式}

\begin{fillblank}[数理人試験 2019 問題1 設問2]
  \noindent \textbf{◯DB則}
  \begin{lawstructure}
    \begin{lawarticle}{第五十八条}{積立不足に伴い拠出すべき掛金の額}
      法第六十三条の厚生労働省令で定めるところにより算定した額は、
      次条第一項前段の規定により（　　A　　）の掛金の額に追加して拠出する場合にあっては第一号の額以上第二号の額以下の範囲内で規約で定める額とする。
    
      \lawitem{一}
      次の表の上欄に掲げる当該事業年度の末日における積立比率
      （積立金の額の最低積立基準額（法第五十八条第二項及び
      法第六十二条に規定する場合に当該事業年度の末日までを計算基準日として（　　B　　）の再計算を行ったときは、当該再計算に基づく最低積立基準額に相当する額（当該再計算に係る給付を法第六十条第三項に規定する給付として同項の規定の例により計算した額をいう。）とする。以下この条及び第六十二条において同じ。）に対する比率をいう。以下この項及び次条において同じ。）の区分に応じて同表の下欄に定める額

      （表略）

      \lawitem{二}
      （略）

      \lawparagraph{2}
      前項の規定は、次条第一項前段の規定により（　　C　　）の掛金の額に追加して
      拠出する場合について準用する。この場合において、前項中「（　　A　　）」
      とあるのは「（　　C　　）」と、「積立金の額」とあるのは
      「積立金の額から（　　D　　）における最低積立基準額の見込額から
      （　　E　　）の最低積立基準額を控除した額を控除した額に（　　A　　）における
      （　　F　　）を加算した額（（　　G　　）が見込まれる場合にあっては
      （　　H　　）を控除した額）」と、「この項及び次条」とあるのは「この項」と読み替えるものとする。
      
      （以下略）

    \end{lawarticle}
  \end{lawstructure}
\end{fillblank}


\begin{answerbox}
  \begin{enumerate}[A: ]
    \item 翌事業年度
    \item 掛金の額
    \item 翌々事業年度
    \item 当該事業年度の翌事業年度
    \item 当該事業年度
    \item 積立金の増加見込額
    \item 積立金の額が減少すること
    \item 積立金の減少見込額
  \end{enumerate}
\end{answerbox}

\begin{truefalse}[年金2 H28 問題1(1)]
    下線部の正誤を判定し、誤っている場合は正しい内容に改めなさい。
\\

  確定給付企業年金において積立比率に応じた必要な掛金を設定する方法により決定した特例
  掛金は、非継続基準の財政検証に抵触した事業年度の\underline{翌々事業年度}の規約に定める時期に拠
  出する。
\end{truefalse}

\begin{answerbox}
  \textbf{誤}

  翌事業年度または翌々事業年度
\end{answerbox}

\begin{shortanswer}[年金2 2025 問題2(3)]
  ある事業年度で非継続基準の財政検証を行った結果、当該事業年度の末日における積立比率 
  (積立金の額の最低積立基準額に対する比率をいう。以下同じ)が 1.0 未満であった。この
  場 合の積立不足に伴い拠出すべき掛金の額について、以下の(A)、(B)に答えなさい。
  (当該額を、翌事業年度の掛金の額に追加して拠出することを規約に定めているものとする)

  \begin{enumerate}[(A)]
    \item 積立比率が 1.0 未満であっても追加の掛金を拠出しないものとすることができる場合が確定給付企業年金法施行規則第 59 条に定められている。その内容を簡潔に入力しなさい。(200 字以内)
    \\
    \item 確定給付企業年金法施行規則第 58 条では、翌事業年度の掛金の額に追加して拠出する額の上限と下限が定められている。1.から3.の積立比率に応じた下限額の算定方法をそれぞれ簡潔に入力しなさい。(各 100 字以内)
    \\
    \begin{enumerate}[1. ]
      \item 積立比率が 0.9 以上 1.0 未満であった場合
      \item 積立比率が 0.8 以上 0.9 未満であった場合
      \item 積立比率が 0.8 未満であった場合
    \end{enumerate}
  \end{enumerate}

\end{shortanswer}

\begin{answerbox}
  \begin{enumerate}[(A)]
    \item 当該事業年度の末日における積立比率が 0.9 以上であって、かつ、当該事業年度の前 3 事業年度のうち少なくとも 2 事業年度の積立比率が 1.0 以上である場合
    \\
    \item \begin{enumerate}[1. ]
      \item 積立金の額が最低積立基準額を下回る額(以下、不足額という)に 15 分の 1 を乗じて得た額
      \item 不足額から最低積立基準額に 0.1 を乗じて得た額を控除した額を 10 で除して得た額に、最低積立基準額に 150 分の 1 を乗じて得た額を加算した額
      \item 不足額から最低積立基準額に 0.2 を乗じて得た額を控除した額を 5 で除して得た額に、最低積立基準額に 60 分の 1 を乗じて得た額を加算した額
    \end{enumerate}
  \end{enumerate}
\end{answerbox}

\newpage

\subsubsection{回復計画方式}

\begin{fillblank}[数理人試験 2025 問題1 設問1 H]
    \noindent \textbf{◯法令解釈}
  \begin{lawstructure}
    \begin{lawarticle}{}{}
      第５ 積立金の積立てに関する事項

      （略）

      \lawparagraph{2}
      確定給付企業年金法施行規則の一部を改正する省令（平成２４年厚生労働省令第１３号）附則第４条
      第１項第３号の「前事業年度における掛金の水準の（略）を上回らないように定められるものである
      こと。」とは、毎事業年度において後の事業年度になるほど掛金の（　　A　　）を大きく見込むな
      ど後の事業年度の負担が過大なものとならないようにすること。
    \end{lawarticle}
  \end{lawstructure}
\end{fillblank}

\begin{answerbox}
  \begin{enumerate}[A: ]
    \item 引き上げ幅
  \end{enumerate}
\end{answerbox}

\begin{truefalse}[年金2 2019 問題1(1)\textcircled{4}]
  下線部の正誤を判定し、誤っている場合は正しい内容に改めなさい。
  \\

  確定給付企業年金において、回復計画における毎事業年度の掛金の額の見込額は、\underline{直近三事
業年度における加入者数}\\
\underline{の実績}を用いて平準的に求められるもの、または前事業年度におけ
る掛金の水準の伸びを上回らないように定められている。
\end{truefalse}

\begin{answerbox}
  \textbf{誤}

  直近五事業年度における加入者数の実績
\end{answerbox}

\begin{truefalse}[年金2 H28 問題1(1)\textcircled{2}]
  下線部の正誤を判定し、誤っている場合は正しい内容に改めなさい。
  \\

  確定給付企業年金において積立比率の回復計画を策定する場合、積立金の額の見込額の計算
  に用いる運用利回りは、財政検証の基準日の属する事業年度の末日における最低積立基準額
  の算定に用いる予定利率、当該事業年度の翌事業年度末日における最低積立基準額の算定に
  用いる予定利率または当該事業年度を含む直近5事業年度における積立金に係る運用利回り
  の実績の平均もしくは掛金額の計算に用いる予定利率のうちいずれか低い率のうち
  \underline{最も高い率}を上回らないものとする必要がある。
\end{truefalse}

\begin{answerbox}
  \textbf{正}
\end{answerbox}

\begin{shortanswer}[年金2 2023 問題2(5)]
  確定給付企業年金制度において非継続基準の財政検証に抵触した場合の積立不足に伴っ
て拠出する掛金額・拠出方法のうち「積立比率の回復計画を作成して積立不足を解消する方
法（確定給付企業年金法施行規則の一部を改正する省令（平成 24 年厚生労働省令第 13 号）
附則第４条）」について、財政検証を行った事業年度の翌々事業年度以後の積立金の額の運用
利回りの見込みに求められる要件を簡潔に入力しなさい。（200字以内）
\end{shortanswer}

\begin{answerbox}
  以下のうち最も高い率を上回らないこと。
  \begin{itemize}
    \item 当該事業年度の末日における最低積立基準額の算定に用いる予定利率
    \item 当該事業年度の翌事業年度の末日における最低積立基準額の算定に用いる予定利率
    \item 当該事業年度を含む直近五事業年度における積立金に係る運用利回りの実績の平均
    もしくは掛金の算定に用いた予定利率のうちいずれか低い率
  \end{itemize}
\end{answerbox}

\newpage

\subsubsection{総合問題}

\begin{shortanswer}[年金2 2020 問題2(1)]
  非継続基準に抵触した場合の積立不足に伴って拠出する掛金の設定方法として、「積立比率
  に応じて必要な掛金を設定する方法」と「回復計画を作成して積立不足を解消する方法」があ
  り、これらは原則として継続して用いることが望ましいとされているが、設定方法を変更でき
  る合理的な理由について具体的事例を 3 つ簡記しなさい。
\end{shortanswer}

\begin{answerbox}
  次のうち３つを記載。
  \begin{itemize}
    \item 合併、分割などにより積立金、最低積立基準額が大幅に変動したとき
    \item 積立金の額の評価方法を変更できる事象が生じたとき
    \item 給付内容の大幅な変更をしたとき
    \item 最低保全給付に係る規約の大幅な変更をしたとき
    \item 平成30年6月22日付規則改正以後、初めて非継続基準に抵触したとき
  \end{itemize}
\end{answerbox}

\newpage

\subsection{積立超過}

\begin{fillblank}[数理人試験 2024 問題1 設問1 A〜F]
    \noindent \textbf{◯DB法}
  \begin{lawstructure}
    \begin{lawarticle}{第六十四条}{積立上限額を超える場合の掛金の控除}
      事業主等は、毎事業年度の決算において、積立金の額が次項に規定する積立上限額を上回っている場合には、当該上回った額を基準として厚生労働省令で定めるところにより算定した額を、第五十五条第三項に定めるところにより算定した掛金の額から厚生労働省令で定めるところにより控除しなければならない。この場合において、当該控除すべき額が同項に定めるところにより算定した掛金の額以上となったときは、当該事業主等に係る掛金については、同条第一項の規定は、適用しない。
    
      \lawparagraph{2}
      積立上限額は、当該確定給付企業年金の（　　A　　）を長期間にわたって確実に確保することができる積立金の水準を上回る額として、厚生労働省令で定めるところにより算定するものとする。
    \end{lawarticle}
  \end{lawstructure}

    \noindent \textbf{◯DB則}
  \begin{lawstructure}
    \begin{lawarticle}{第六十条}{積立上限額を超える場合の掛金の控除額}
      法第六十四条第一項の厚生労働省令で定めるところにより算定した額は、次のいずれかの額とする。

      \lawitem{一}
      当該事業年度の末日において積立金の額が法第六十四条第二項に規定する積立上限額
      （以下「積立上限額」という。）を上回った額のうち未だ控除していない額に、
      当該未だ控除していない額に係る当該事業年度の末日から（　　B　　）までの期
      間に応ずる利子に相当する額（以下この条において「利子相当額」という。）を加
      算した額又は控除前の掛金の額のいずれか（　　C　　）

      \lawitem{二}
      次条第一号の控除を開始するときから（　　D　　）までの期間において、
      積立金の額が積立上限額を上回った額と当該上回った額に係る利子相当額の合
      計額を掛金の額から均等に控除する場合の額又は控除前の掛金の額のいずれか（　　C　　）

      \lawparagraph{2}
        前項の利子相当額の計算に用いる利率は、（　　E　　）における下限予定利率とする。
    \end{lawarticle}
  \end{lawstructure}

  \begin{lawstructure}
    \begin{lawarticle}{第六十一条}{掛金の控除の方法}
      法第六十四条第一項の掛金の額からの控除は、規約で定めるところにより、前条の規定により算定した額を次のとおり控除するものとする。

      \lawitem{一}
      遅くとも当該事業年度の翌々事業年度の最初に拠出する掛金の額から控除を開始すること。

      \lawitem{二}
      掛金の一部を加入者が負担している場合にあっては、
      当該掛金の額からの控除後に加入者が負担する掛金の額が当該加入者に係る当該掛金
      の額からの（　　F　　）を超えないこと。
    \end{lawarticle}
  \end{lawstructure}
  
\end{fillblank}

\begin{answerbox}
  \begin{enumerate}[A: ]
    \item 財政の安定性
    \item 控除する日
    \item 小さい額
    \item 当該事業年度の翌々事業年度の末日
    \item 当該事業年度の末日
    \item 控除後の掛金の額の二分の一
  \end{enumerate}
\end{answerbox}



\begin{fillblank}[数理人試験 2022 問題4 設問1　\&　数理人試験 2024 問題1 設問1 E〜H]
    \noindent \textbf{◯DB則}
  \begin{lawstructure}
    \begin{lawarticle}{第六十二条}{積立上限額の算定方法}
      当該事業年度の末日における積立上限額は、次のいずれか大きい額に一・五を乗じて得た額とする。

      \lawitem{一}
      次の要件を満たす基礎率を用いて計算したにおける数理債務の額

      \hangedlawitem{イ}{3}
      予定利率は、（　　A　　）における下限予定利率とすること。

      \hangedlawitem{ロ}{3}
      予定死亡率は、基準死亡率に、次に掲げる加入者、加入者であった者又はその遺族等の区分に応じそれぞれ定める率を乗じた率とすること。

      \hangedlawitem{（１）}{4}
      加入者　（　　B　　）

      \hangedlawitem{（２）}{4}
      男子であって、加入者であった者又はその遺族（（４）に掲げる者を除く。）　（　　C　　）

      \hangedlawitem{（３）}{4}
      女子であって、加入者であった者又はその遺族（（４）に掲げる者を除く。）　（　　D　　）

      \hangedlawitem{（４）}{4}
      障害給付金の受給権者　（　　E　　）（（１）に掲げる者を除く。）

      \hangedlawitem{ハ}{3}
      その他の基礎率は、（　　F　　）とすること。

      \lawitem{二}
      （　　G　　）
    \end{lawarticle}
  \end{lawstructure}
\end{fillblank}

\begin{answerbox}
  \begin{enumerate}[A: ]
    \item 当該事業年度の末日
    \item 零
    \item 〇・六八
    \item 〇・六五
    \item 一・〇
    \item 前回の財政計算で使用した基礎率
    \item 当該事業年度の最低積立基準額
  \end{enumerate}
\end{answerbox}


\begin{fillblank}[数理人試験 H17 問題1(4)　\&　数理人試験 H26 問題1(4)]
    \noindent \textbf{◯DB則}
  \begin{lawstructure}
    \begin{lawarticle}{第六十二条}{積立上限額の算定方法}
      当該事業年度の末日における積立上限額は、次のいずれか大きい額に（　　A　　）を乗じて得た額とする。

      \lawitem{一}
      次の要件を満たす基礎率を用いて計算したにおける（　　B　　）の額

      \hangedlawitem{イ}{3}
      予定利率は、当該事業年度の末日における（　　C　　）とすること。

      \hangedlawitem{ロ}{3}
      予定死亡率は、（　　D　　）に、次に掲げる加入者、加入者であった者又はその遺族等の区分に応じそれぞれ定める率を乗じた率とすること。

      \hangedlawitem{（１）}{4}
      加入者　（　　E　　）

      \hangedlawitem{（２）}{4}
      男子であって、加入者であった者又はその遺族（（４）に掲げる者を除く。）　（　　F　　）

      \hangedlawitem{（３）}{4}
      女子であって、加入者であった者又はその遺族（（４）に掲げる者を除く。）　（　　G　　）

      \hangedlawitem{（４）}{4}
      障害給付金の受給権者　（　　H　　）（（１）に掲げる者を除く。）

      \hangedlawitem{ハ}{3}
      その他の基礎率は、（　　I　　）とすること。

      \lawitem{二}
      （　　J　　）
    \end{lawarticle}
  \end{lawstructure}
\end{fillblank}

\begin{answerbox}
  \begin{enumerate}[A: ]
    \item 一・五
    \item 数理債務
    \item 下限予定利率
    \item 基準死亡率
    \item 零
    \item 〇・六八
    \item 〇・六五
    \item 一・〇
    \item 前回の財政計算で使用した基礎率
    \item 当該事業年度の最低積立基準額
  \end{enumerate}
\end{answerbox}




\begin{truefalse}[年金2 2025 問題1(1)ア\ajRoman{3}]
  以下の文章の正誤を判定し、誤っている場合は正しい内容に改めなさい。
  \\

  事業主等は、毎事業年度の決算において、積立金の額が積立上限額を上回ってい
  る場合には、当該上回った額を基準として算定した額を掛金の額から控除しなけ
  ればならない。この積立金の額は、掛金の額の計算に用いる積立金の額の評価方
  法にかかわらず時価で評価するものとする。
\end{truefalse}

\begin{answerbox}
  \textbf{誤}

  掛金の額の計算に用いる積立金の額の評価方法を用いる。
\end{answerbox}

\begin{truefalse}[年金2 2025 問題1(1)ウ\ajRoman{3}]
    以下の文章の正誤を判定し、誤っている場合は正しい内容に改めなさい。
  \\

  積立上限額を算定するための数理債務を算定する場合において基準死亡率に乗じ
  る率は、加入者について、男子にあっては 0.68、女子にあっては 0.65 である。
\end{truefalse}

\begin{answerbox}
  \textbf{誤}

  ×加入者→◯加入者であった者またはその遺族（障害給付金の受給権者を除く）
\end{answerbox}

\begin{shortanswer}[年金2 H29 問題2(2)\textcircled{1}]
  数理債務、最低積立基準額および積立上限額の算定に用いる数理債務の計算に用いる予定死
亡率の決め方を簡記しなさい
\end{shortanswer}

\begin{answerbox}
  ＜数理債務＞

  加入者等及びその遺族の性別及び年齢に応じた死亡率として厚生労働大臣が定める率（以下、「基準死亡率」という）とする。
ただし、次の範囲内で定めた率を基準死亡率に乗じたものとすることができる。

加入者：0以上

加入者であった者またはその遺族（男子）：0.68以上1.0以下

加入者であった者またはその遺族（女子）：0.65以上1.0以下

障害給付金の受給権者：1.0以上
\\

＜最低積立基準額＞

基準死亡率に以下の率を乗じたものとする。

男子は0.84、女子は0.825。
\\

＜積立上限額＞

基準死亡率に以下の率を乗じたものとする。

加入者：0

加入者であった者またはその遺族（男子）：0.65

加入者であった者またはその遺族（女子）：0.68

障害給付金の受給権者：1.0

\end{answerbox}

\begin{shortanswer}[年金2 2024 問題2(3)]
  事業主等は、確定給付企業年金の毎事業年度の決算において、積立金の額が積立上限額を上
  回っている場合には、当該上回った額を基準として「厚生労働省令で定めるところにより算
  定した額」を、掛金の額から控除しなければならない。これについて、以下の(A)、(B)
  に答えなさい。
\\
  \begin{enumerate}[(A)]
    \item 掛金の額から控除を開始する時期として確定給付企業年金法施行規則第61条に規定する内容を簡潔に入力しなさい。(100 字以内)
    \\
    \item 「厚生労働省令で定めるところにより算定した額」として確定給付企業年金法施行規則第60条に規定する内容を 2 つ簡潔に入力しなさい。(400 字以内)
  \end{enumerate}
\end{shortanswer}

\begin{answerbox}
  \begin{enumerate}[(A)]
    \item 規約で定めるところにより、遅くとも当該事業年度の翌々事業年度の最初に拠出する掛金の額から控除を開始する。
    \\
    \item 次のいずれかの額とする。\\
＜前詰め控除方式＞\\
当該事業年度の末日において積立金の額が積立上限額を上回った額のうち未だ控除していない額に、当該未だ控除していない額に係る当該事業年度の末日から控除する日までの期間に応ずる加減予定利率による利子に相当する額（以下、「利子相当額」という）を加算した額または控除前の掛金の額のいずれか小さい額
\\
＜均等控除方式＞\\
控除を開始するときから当該事業年度の翌々事業年度の末日までの期間において、積立金の額が積立上限額を上回った額と当該上回った額に係る利子相当額の合計額を掛金の額から均等に控除する場合の額または控除前の掛金の額のいずれか小さい額
  \end{enumerate}
  
\end{answerbox}

\begin{explanationbox}
  類題：数理人試験 2022 問題4 設問2
\end{explanationbox}

\begin{shortanswer}[数理人試験 H19 問題4]
  ある確定給付企業年金制度の事業年度末の決算において、積立金の額が積立上限額を上回っ
  た。この場合の掛金の控除額及び控除方法について以下の設問に答えよ。
  なお、掛金の控除方法は、確定給付企業年金法施行規則(以下「規則」という。)第 60 条第 1 項第 1 号
  による方法及び同項第 2 号による方法の2つがあるが、本問において、
  それぞれ 「前詰め控除方式」、「均等控除方式」という。
  \begin{enumerate}[(1)]
    \item 掛金の控除を行う場合に、次のAからCに関する条件を述べよ。 
    \begin{enumerate}[A. ]
      \item 掛金の控除開始時期
      \item 掛金の加入者負担額
      \item 均等控除方式の場合における掛金控除終了時期
    \end{enumerate}
    \item 当該事業年度末の決算(財政検証)時における各数値が以下のとおりであった。
    この場合、 「前詰め控除方式」、「均等控除方式」のそれぞれの方法による控除後掛金月額(平成 20 年 4 月分)
    を計算せよ。なお、計算の過程を必ず解答用紙に示すこと。(解答にあたっては、 計算過程及び最終計算結果について、
    小数第 2 位を四捨五入すること。)

    \begin{table}[htbp]
      \centering
      \begin{tabular}{|l|}\hline
        事業年度の末日(財政検証基準日):平成 19 年 3 月 31 日 \\
        財政計算上の予定利率:4.5\% \\
        純資産額:10,000 \\
        数理上資産額:9,500 \\
        規則第 62 条第 1 号に定める要件を満たす基礎率を用いて\\
          \;\;\;\; 計算した(積立上限額の算定に使用する)数理債務:6,000\\
        最低積立基準額:5,000\\
        掛金拠出時期:毎月(年 12 回)\\
        掛金控除開始時期:平成 20 年 4 月\\
        控除前の掛金月額(加入者負担なし):600\\
        確定年金現価率(予定利率 4.5\%、控除期間 1 年、年 12 回拠出):11.7\\
        確定年金現価率(予定利率 4.5\%、控除期間 2 年、年 12 回拠出):22.9\\
        確定年金現価率(下限予定利率 = 1.2\%、控除期間 1 年、年 12 回拠出):11.9 \\
        確定年金現価率(下限予定利率 = 1.2\%、控除期間 2 年、年 12 回拠出):23.7 \\ \hline
      \end{tabular}
    \end{table}
  \end{enumerate}
\end{shortanswer}

\begin{answerbox}
  \begin{enumerate}[(1)]
    \item \begin{enumerate}[A.]
      \item 遅くとも当該事業年度の翌々事業年度の最初に拠出する掛金の額から控除を開始すること
      \item 掛金の一部を加入者が負担している場合には、掛金控除後の加入者負担額が、掛金控除後の掛金額の 2 分の 1 を超えないこと
      \item (掛金控除開始時から)当該事業年度の翌々事業年度の末日までの期間において掛金を控除すること
    \end{enumerate}
    \item \begin{flalign*}
        積立上限額 &= \max(数理債務, 最低積立基準額) \times 1.5 \\
          &=  \max(6,000, 5,000) \times 1.5 = 9,000
      \end{flalign*}
      
      平成19年3末時点で、数理上資産額が積立上限額を上回った額：$9,500-9,000 = 500$

      平成19年3末時点の加減予定利率が1.2\%なので、

      平成20年4月からの控除対象額：$500 \times 1.012 = 506$

      従って、控除後掛金月額は、

      【前詰め控除方式】

      平成 20 年 4 月分の控除後掛金月額:$600-506=94.0$

      【均等控除方式】

      平成 20 年 4 月〜平成 21 年 3 月までの期間、均等に掛金控除するため、 
      
      控除額(月額)：$506/11.9 = 42.521\cdots = 42.5$

      平成 20 年 4 月分の控除後掛金月額：$600-42.5=557.5$
  \end{enumerate}
\end{answerbox}

\newpage

\subsection{臨時掛金}

出題例なし

\newpage

\subsection{簡易な基準に基づくDBの財政検証特例}

出題例なし

\newpage

\subsection{規約型DBの積立金の管理運用}

出題例なし

\newpage

\subsection{基金型DBの積立金の管理運用}

出題例なし

\newpage

\subsection{分散投資と政策的資産構成割合の策定}


\begin{InyouJobunbox}
  \begin{itemize}
    \item 法第67条
    \item 令第46条〜第46条の2
    \item 則第84条、第84条の4〜6
    \item 解釈第6-3、6-5
  \end{itemize}
\end{InyouJobunbox}

\begin{truefalse}[年金1 2021 問題1(3)\textcircled{1}]
  下線部の正誤を判定し、誤っている場合は正しい内容に改めなさい。
\\

  2018 年（平成30 年）4 月より、積立金の額が100 億円以上の確定給付企業年金制度におい
て資産運用委員会の設置が義務づけられているが、2020 年（令和2 年）10 月1 日以後最初に
終了する事業年度の翌事業年度からその設置要件が明確化され、直近3 年の財政決算のうち少
なくとも2 年で、 \underline{「積立金の額が100 億円以上」かつ「責任準備金の額または最低積立基準額の
いずれか高い額が100 億円以上」}に該当する確定給付企業年金制度とされている。
\end{truefalse}

\begin{answerbox}
  \textbf{誤}

  「積立金の額が100億円以上」または「責任準備金の額または最低積立基準額のいずれか低い額が100億円以上」
\end{answerbox}

\begin{truefalse}[年金2 2023 問題1(1)ウ]
  以下の文中の(\ajRoman{1})〜(\ajRoman{4})のうち、適当なものを全て選びなさい。
\\

  確定給付企業年金において資産運用委員会の設置が求められる具体的な要件は、次の
\fbox{（\ajRoman{1}）\textcircled{1}と\textcircled{2}のいずれにも}該当することである。

\textcircled{1} 直近３カ年の決算のうち少なくとも\fbox{（\ajRoman{2}）２カ年}において、積立金（純資産）の額が
\fbox{（\ajRoman{3}）２００億円}以上であること。

\textcircled{2}  直近３カ年の決算のうち少なくとも\fbox{（\ajRoman{2}）２カ年}において、\fbox{（\ajRoman{4}）数理債務}の額又は最
低積立基準額のいずれか低い額が\fbox{（\ajRoman{3}）２００億円}以上であること。
\end{truefalse}

\begin{answerbox}
  \ajRoman{2}
\end{answerbox}

\begin{explanationbox}
  \ajRoman{1}：\textcircled{1}と\textcircled{2}のいずれか

  \ajRoman{2}：正

  \ajRoman{3}：100億円

  \ajRoman{4}：責任準備金
\end{explanationbox}

\newpage

\subsection{運用の基本方針}


\begin{InyouJobunbox}
  \begin{itemize}
    \item 令第45条
    \item 則第82条〜第83条、第84条の2〜3
    \item 解釈第6-1〜6-2
  \end{itemize}
\end{InyouJobunbox}

\begin{fillblank}[年金2 2019 問題1(4)]
  \noindent \textbf{◯DB令}
  \begin{lawstructure}
    \begin{lawarticle}{第四十五条}{運用の基本方針}
      事業主（厚生労働省令で定める要件に該当する規約型企業年金を実施するものを除く。以下この条において同じ。）及び基金は、積立金の運用に関して、運用の目的その他厚生労働省令で定める事項を記載した基本方針（以下この条及び第四十六条の二第三項において「基本方針」という。）を作成し、当該基本方針に沿って運用しなければならない。
    
      （以下略）
    \end{lawarticle}
  \end{lawstructure}

  \noindent \textbf{◯DB則}
  \begin{lawstructure}
    \begin{lawarticle}{第八十二条}{基本方針を定めることを要しない規約型企業年金の要件}
      令第四十五条第一項の厚生労働省令で定める要件は、当該確定給付企業年金が（　　A　　）であることとする。
    \end{lawarticle}
  \end{lawstructure}

  \begin{lawstructure}
    \begin{lawarticle}{第八十四条の二}{運用の基本方針の作成又は変更に当たって加入者の意見を聴く方法}
      令第四十五条第三項（同条第五項において準用する場合を含む。）の規定により加入者の意見を聴く場合には、次のいずれかの方法により行うものとする。

      \lawitem{一}
      規約で定めるところにより加入者の代表者を選任し、必要に応じて当該代表者が参画する委員会を設置して次に掲げる措置を講ずる方法

      \hangedlawitem{イ}{3}
      基本方針を作成又は変更する際に、当該代表者に意見を述べる機会を与えること。

      \hangedlawitem{ロ}{3}
      年一回以上、基本方針に関して、当該代表者に意見を述べる機会を与えること。

      \hangedlawitem{ハ}{3}
      当該代表者からの求めがあった場合に、毎事業年度の積立金の資産の額その他（　　B　　）を
      当該代表者に開示すること。

      \lawitem{二}
      基金型企業年金にあっては、次に掲げる措置を講ずる方法

      \hangedlawitem{イ}{3}
      基本方針を作成又は変更する際に、規約で定めるところにより加入者に意見の提出の機会を与えること。

      \hangedlawitem{ロ}{3}
      基本方針を作成又は変更する際に、（　　C　　）を経ること。

      \hangedlawitem{ハ}{3}
      代議員からの求めがあった場合に、毎事業年度の積立金の資産の額その他（　　B　　）を当該代議員に開示すること。

      \lawitem{三}
      次に掲げる確定給付企業年金以外の確定給付企業年金にあっては、第八十七条の規定に基づき周知される基本方針に関して意見を聴く方法

      \hangedlawitem{イ}{3}
      第二十九条第三号の積立金の運用利回りの実績に基づき令第二十四条第一項第三号の再評価若しくは
      同条第三項の改定を行う確定給付企業年金（第二十九条第四号又は第五号において同条第三号の積立金
      の運用利回りの実績を用いるものを含み、国債、保険業法施行規則第七十五条の二第一項第一号に規定する
      一般勘定を設ける保険契約に係る資産その他これらに準ずる資産のみで資産を構成し、資産の構成割合を
      あらかじめ規約で定めるもの及び（　　A　　）を除く。）

      \hangedlawitem{ロ}{3}
      （　　D　　）

      \lawparagraph{2}
      前項第一号の加入者の代表者は、規約で定めるところにより、（　　E　　）を有する代理人に同号イ及びロの意見を述べさせることができる。

      \lawparagraph{3}
      第一項第三号イ又はロに掲げる確定給付企業年金を実施する事業主又は基金は、基本方針の作成又は変更に当たって、第一項第一号イ若しくはロ又は第二号イの意見を十分に考慮しなければならない。
    \end{lawarticle}
  \end{lawstructure}
\end{fillblank}


\begin{answerbox}
  \begin{enumerate}[A: ]
    \item 受託保証型確定給付企業年金
    \item 積立金の運用の実績
    \item 代議員会の議決
    \item リスク分担型企業年金
    \item 専門的知識及び経験
  \end{enumerate}
\end{answerbox}

\begin{fillblank}[年金1 2023 問題1(1)ウ、エ]
  \noindent \textbf{◯法令解釈}
  \begin{lawstructure}
    \begin{lawarticle}{}{}
      第６　積立金の運用に関する事項

      １　運用の基本方針について

      （略）

      なお、運用の基本方針の作成又は変更に当たっては、令第４５条第３項の規定に基
      づき、加入者の意見を聴くこと。

      当該加入者の意見を聴く方法は、規則第８４条の２第１項に規定する方法により実
      施すること。なお、加入者の意見を聴く方法については次の\textcircled{1}から\textcircled{5}までに留意する
      こと。

      \lawitem{\textcircled{1}}
      規則第８４条の２第１項第１号の方法により行う場合で、当該加入者の代表者
      が参画する委員会を設置する場合にあっては、同条第２項の規定に基づき、専門
      的知識及び経験を有する（　　A　　）を当該委員会に参加させることも可能であるこ
      と。


      \lawitem{\textcircled{2}〜\textcircled{4} （略）}

      \lawitem{\textcircled{5}}
      リスク分担型企業年金の場合は、規則第８４条の２第１項第１号ハ及び第２号
      ハに基づく情報開示と同様に、調整率（規則第２５条の２第１項に規定する調整
      率をいう。以下同じ。）及び（　　B　　）（規則第２５条の２第２項に規定する（　　B　　）
      をいう。以下同じ。）に係る情報について、加入者の代表者又は代議員から
      の求めがあった場合に開示すること。

      （以下略）
    \end{lawarticle}
  \end{lawstructure}
\end{fillblank}

\begin{answerbox}
  \begin{enumerate}[A: ]
    \item 代理人
    \item 超過比率
  \end{enumerate}
\end{answerbox}

\begin{truefalse}[年金1 2020 問題1(3)\textcircled{1}]
  下線部の正誤を判定し、誤っている場合は正しい内容に改めなさい。
  \\

  規約型企業年金（ただし、受託保証型確定給付企業年金を除く）において事業主が積立金の
運用に関する基本方針を作成しようとするときは加入者の意見を聴かなければならないが、具
体的に講じるべき措置の1 つとして「加入者の代表者からの求めがあった場合に、\underline{毎事業年度
の責任準備金の額および積立金の資産の額}を当該代表者に開示すること」がある。
\end{truefalse}

\begin{answerbox}
  \textbf{誤}

  毎事業年度の積立金の資産の額その他積立金の運用の実績
\end{answerbox}

\begin{truefalse}[年金1 H29 問題2(1)\textcircled{4}]
  以下の文章の正誤を判定し、誤っている場合は正しい内容に改めなさい。
  \\

  規約型のリスク分担型企業年金（運用の基本方針を策定済み。）において、事業主が運用の基
本方針について加入者の代表者に意見を述べる機会を与える必要があるのは、基本方針を変更
する際のみである。
\end{truefalse}

\begin{answerbox}
  \textbf{誤}

  基本方針を変更する際だけでなく、年１回以上、意見を述べる機会を与える必要がある。
\end{answerbox}

\newpage

\subsection{政令への委任}

出題例なし

\newpage